% Paper A. Must stand without any fluid-universe material.
% Every number here regenerates from `python -m phenomenology.run_all`.
\documentclass[11pt]{article}
\usepackage{amsmath,amssymb,booktabs,hyperref}
\title{A $Z_3$ circulant reparametrisation of the charged-lepton masses:\\
algebra, pole-mass fit, and one-loop running}
\author{Research note generated from the \texttt{fluid-universe} repository}
\date{\today}

\begin{document}
\maketitle

\begin{abstract}
We test the claim that a $Z_3$ multiplet structure forces a Hermitian circulant
matrix whose eigenvalues are the square roots of the charged-lepton masses.
The algebra is correct and, in this form, is already present in the unpublished
circulant treatment of Brannen (2006): a Hermitian circulant with three real
parameters $(a,c,\delta)$ has eigenvalues
$\lambda_k=a+2c\cos(\delta+2\pi k/3)$, the Koide ratio is
$Q=1/3+2c^2/(3a^2)$, and $Q=2/3$ is equivalent to $c/a=1/\sqrt2$.
The construction does \emph{not} predict $Q=2/3$: three parameters reproduce
three masses exactly, and Koide is one additional condition. The natural
quantum-rotor route is excluded: with a real $\cos3\phi$ potential the $k=1$
and $k=2$ sectors are exactly degenerate, so the naive map yields only two
distinct roots. Fitting PDG pole
masses gives $Q=0.66666446$ and $c/a=0.70710444$. The exact statement that
$c/a$ is one-loop renormalisation-group invariant is false; the
generation-dependent $y_k^2$ term breaks it. For charged leptons the numerical
drift is only $\Delta Q=+2.89\times10^{-6}$, $\Delta r=+3.06\times10^{-6}$ from
$M_Z$ to $10^{16}$ GeV. The full-SM $\overline{\mathrm{MS}}$ masses at $M_Z$
give $Q=0.6678240$, and the one-pass QED conversion of the pole masses alone
gives $0.6679160$; both are about $0.18\%$ above $2/3$. These are
scheme/matching effects, not RG running. We find no new falsifiable prediction beyond the built-in Koide
value. The $Z_3$ origin and the value $c/a=1/\sqrt2$ remain assumptions.
\end{abstract}

\section{Setup and conventions}
\label{sec:setup}

Let $\omega=e^{2\pi i/3}$ and
\begin{equation}
M=\mathrm{circ}(a,b,b^*)
=\begin{pmatrix}
a & b & b^*\\
b^* & a & b\\
b & b^* & a
\end{pmatrix},
\qquad b=c\,e^{i\delta},
\label{eq:M}
\end{equation}
with $a,c,\delta\in\mathbb{R}$ and $c\ge0$. \emph{Derived:} $M=M^\dagger$; for
generic $\delta$, $M\ne M^T$, so the matrix is Hermitian but not complex
symmetric. \emph{Assumed:} the three charged leptons form a $Z_3$ multiplet and
$M$ acts on the vector of signed square roots
$x=(\sqrt{m_e},\sqrt{m_\mu},\sqrt{m_\tau})$. The second assumption is the
content of the hypothesis; nothing below derives it.

\section{Algebra}
\label{sec:algebra}

\emph{Derived.} The eigenvectors of \eqref{eq:M} are the DFT columns
$v_k=(1,\omega^k,\omega^{2k})/\sqrt3$, independent of $\delta$, and
\begin{equation}
\lambda_k=a+2c\cos\!\left(\delta+\frac{2\pi k}{3}\right),\qquad k=0,1,2.
\label{eq:eig}
\end{equation}
Consequently
\begin{equation}
\sum_k\lambda_k=3a,\qquad
\sum_k\lambda_k^2=3a^2+6c^2,
\end{equation}
and, for the signed roots,
\begin{equation}
Q\equiv\frac{\sum_k\lambda_k^2}{\left(\sum_k\lambda_k\right)^2}
=\frac13+\frac{2c^2}{3a^2}.
\label{eq:Q}
\end{equation}
Equation \eqref{eq:Q} is verified symbolically and by $10^4$ random numeric
draws; $Q=1/3$ at $c=0$, $Q=2/3$ at $c^2/a^2=1/2$, and $Q=1$ at $c^2/a^2=1$.
Hence
\begin{equation}
Q=\frac23 \iff \frac{c}{a}=\frac{1}{\sqrt2}.
\label{eq:koide}
\end{equation}

\paragraph{Sign caveat.} Physical masses use $m_k=\lambda_k^2$ and
$Q_{\rm phys}=\sum m_k/(\sum\sqrt{m_k})^2=\sum\lambda_k^2/(\sum|\lambda_k|)^2$.
Equation \eqref{eq:Q} holds for $Q_{\rm phys}$ only when every $\lambda_k\ge0$.
Counterexample: $a=1$, $c=0.9$, $\delta=\pi$ gives signed roots
$(-0.8,1.9,1.9)$, $Q_{\rm signed}=0.8733$ but $Q_{\rm phys}=0.3715$. The lepton
fit below satisfies $\lambda_k>0$, so the caveat is inactive there but it
forbids stating \eqref{eq:Q} as a universal mass identity.

\paragraph{No prediction.} Equation \eqref{eq:koide} is a condition on a free
ratio. The map $(a,c,\delta)\mapsto(\sqrt{m_e},\sqrt{m_\mu},\sqrt{m_\tau})$ is
invertible on positive roots, so the fit is exact for \emph{any} three masses.
A $Z_3$ multiplet alone therefore does not explain Koide. The universal bound
$c/a\le1/2$ is false: positive roots $(10,1,1)$ give $c/a=0.75$ with $Q=0.7083$.

\paragraph{The rotor route is excluded.} A natural attempt is to generate the
three roots from a quantum rotor $H=L^2/(2I)+\mu(1-\cos3\phi)$. Exact
diagonalisation in the Fourier sector $n=k+3j$ gives, at $I=\mu=1$,
\begin{equation}
E_0=0.8911861597,\qquad E_1=E_2=1.3170345140,
\label{eq:rotor}
\end{equation}
with the degeneracy exact because $\phi\to-\phi$ maps sector $1$ to sector $2$.
The naive map $a=(\ell_0+2\ell_1)/3$, $c=(\ell_1-\ell_0)/3$ therefore produces
roots $(1.45898,1.03314,1.03314)$: only two distinct values, not three. The
spectrum converges to $10^{-9}$ between cutoff $n_{\max}=20$ and $160$, and the
free-rotor limit reproduces $E_n=n^2/(2I)$. Thus the rotor model, as stated, is
falsified rather than merely unverified. The associated domain-wall tension
$T=(8/3)\sqrt{\mu}$ is verified analytically and numerically, but that identity
alone does not rescue the mass map.

\section{Pole-mass fit}
\label{sec:fit}

\emph{Numeric.} Inputs are the PDG 2024 pole masses
$m_e=0.51099895000(15)$ MeV, $m_\mu=105.6583755(23)$ MeV, and
$m_\tau=1776.93(9)$ MeV. Inverting \eqref{eq:eig} gives
\begin{center}
\begin{tabular}{lr}
\toprule
$a$ & $17.715838478\ \mathrm{MeV}^{1/2}$\\
$c$ & $12.526948122\ \mathrm{MeV}^{1/2}$\\
$c/a$ & $0.707104444$\\
$\delta \bmod 2\pi/3$ & $0.22222476$ rad\\
$Q$ & $0.666664463$\\
$\max|\text{residual}|$ & $1.0\times10^{-14}\ \mathrm{MeV}^{1/2}$\\
\bottomrule
\end{tabular}
\end{center}
A fixed-seed Monte Carlo ($10^5$ draws, PDG $1\sigma$ errors) gives
$\sigma_Q=5.1\times10^{-6}$ and $\sigma_{c/a}=5.4\times10^{-6}$, so
$(Q-2/3)/\sigma_Q=-0.44$. Solving $Q=2/3$ with the measured $m_e,m_\mu$ gives
$m_\tau=1776.969$ MeV, to be compared with $1776.93(9)$ MeV. This is a
post-hoc consistency check of a two-parameter relation, not a prediction of
this paper; it is the known Koide prediction.

The phase has a cyclic ambiguity $\delta\to\delta+2\pi/3$ that leaves the
spectrum invariant. Its minimal representative is numerically close to $2/9$
(difference $2.5\times10^{-6}$ rad). \emph{Speculation:} we attach no
significance to this near-miss. Without a stated prior, rational
approximations of $\delta$ are not evidence.

\section{One-loop running}
\label{sec:rg}

\emph{Derived (standard).} The one-loop Standard-Model RGE for a charged-lepton
Yukawa eigenvalue is
\begin{equation}
16\pi^2\frac{d\ln y_k}{dt}
=\frac32 y_k^2+T-\frac94 g_2^2-\frac{15}{4}g_1^2,
\qquad
T=\mathrm{Tr}\!\left(3Y_u^\dagger Y_u+3Y_d^\dagger Y_d+Y_e^\dagger Y_e\right),
\label{eq:rge}
\end{equation}
with $g_1$ SU(5)-normalised. The gauge and $T$ terms are generation-blind; only
$\frac32y_k^2$ distinguishes the three leptons. Therefore $Q$ and $r=c/a$ are
\emph{not} exactly RG-invariant, contrary to the original hypothesis.

\emph{Numeric.} We integrate \eqref{eq:rge} with the full-SM
$\overline{\mathrm{MS}}$ masses of Huang and Zhou (2021) at $M_Z$,
$m_e=0.48307$ MeV, $m_\mu=101.766$ MeV, $m_\tau=1728.56$ MeV, together with
$m_t(M_Z)=168.26$ GeV and $m_b(M_Z)=2.839$ GeV, and one-loop gauge running.
The result is
\begin{center}
\begin{tabular}{lrr}
\toprule
 & $M_Z$ & $10^{16}$ GeV\\
\midrule
$Q$ & $0.667823994$ & $0.667826882$\\
$c/a$ & $0.708333249$ & $0.708336306$\\
\bottomrule
\end{tabular}
\end{center}
so $\Delta Q=+2.89\times10^{-6}$ and $\Delta r=+3.06\times10^{-6}$ over 32.3
e-folds. The exact-invariance claim is falsified; the numerical stability is
real but is a consequence of the smallness of $y_\tau^2$, not a symmetry.

\paragraph{Scheme versus running.} Converting the pole masses to
$\overline{\mathrm{MS}}$ at $M_Z$ with the published one-pass QED formula
($\alpha(M_Z)=1/128.89$) gives $Q=0.6679160$, in agreement with the value
$0.66792$ of Xing and Zhang (2006). The full-SM matching of Huang and Zhou
instead gives $Q(M_Z)=0.6678240$. Both differ from the pole-mass $Q$ by about
$0.18\%$, an order of magnitude larger than the entire subsequent RG drift.
The scheme/matching conversion and the running must not be conflated.

\paragraph{Truncation.} The absolute high-scale masses are sensitive to a
near-cancellation between the top-Yukawa and gauge terms; at $10^8$ GeV the
one-loop masses differ from the three-loop literature values by up to a few
percent. No claim here uses absolute high-scale masses. A genuine two-loop run
is future work. An earlier estimate of a two-loop drift
$\sim4\times10^{-9}$ per decade had no supporting calculation and is retracted.

\section{What does not survive}
\label{sec:negative}

\begin{itemize}
\item The $Z_3$ multiplet does not force $Q=2/3$; Koide is the extra condition
$c/a=1/\sqrt2$. The value of $c/a$ is not derived.
\item $Q=2/3$ holds for pole masses, not for $\overline{\mathrm{MS}}$ masses at
$M_Z$, where the full-SM matching gives $Q=0.6678240$.
\item $c/a$ is not one-loop RG-invariant.
\item The rotor-to-circulant map cannot give three masses: $E_1=E_2$ exactly.
\item The bound $c/a\le1/2$ is false.
\item No new falsifiable prediction beyond Koide has been found.
\end{itemize}

\section{Status and outlook}
\label{sec:outlook}

The algebra in Sec.~\ref{sec:algebra} is likely known; the circulant
eigenvalue problem and the identification $\lambda_k=\sqrt{m_k}$ appear in
Brannen's unpublished note, and the running-mass behaviour is discussed by
Xing and Zhang and by Sumino. This note contributes a clean separation of three
effects that are often conflated---the exact algebra, the scheme conversion,
and the RG drift---and an explicit statement of what the $Z_3$ hypothesis does
not do.

Any future claim of a dynamical origin must supply a mechanism that fixes
$c/a=1/\sqrt2$ and yields at least one prediction beyond the three charged
lepton masses. Absent that, the construction is a reparametrisation, not a
theory.

\begin{thebibliography}{9}
\bibitem{koide} Y.~Koide, Lett.\ Nuovo Cim.\ \textbf{34} (1982) 201;
Phys.\ Rev.\ D \textbf{28} (1983) 252.
\bibitem{brannen} C.~A.~Brannen, \emph{The Lepton Masses}, unpublished (2006),
\texttt{https://brannenworks.com/MASSES2.pdf}.
\bibitem{xingzhang} Z.-z.~Xing and H.~Zhang,
\emph{On the Koide-like relations for the running masses of charged leptons,
neutrinos and quarks}, arXiv:hep-ph/0602134.
\bibitem{sumino} Y.~Sumino, JHEP \textbf{05} (2009) 075, arXiv:0812.2090.
\bibitem{huangzhou} G.-y.~Huang and S.~Zhou, Phys.\ Rev.\ D \textbf{103} (2021)
016010, arXiv:2009.04851.
\bibitem{xingzhangzhou} Z.-z.~Xing, H.~Zhang, and S.~Zhou,
Phys.\ Rev.\ D \textbf{77} (2008) 113016, arXiv:0712.1419.
\bibitem{arason} H.~Arason \emph{et al.}, Phys.\ Rev.\ D \textbf{46} (1992) 3945.
\bibitem{pdg} S.~Navas \emph{et al.} (Particle Data Group),
Phys.\ Rev.\ D \textbf{110} (2024) 030001.
\end{thebibliography}

\end{document}
